
\subsection*{3. Fibrés projectifs et fibrés en drapeaux}

\textbf{Proposition 3.1.} \uline{Soient} @@M0@@ \uline{un schéma noethérien}, @@M1@@ \uline{un} @@M2@@\uline{-Module localement libre de rang} @@M3@@, @@M4@@ \uline{une suite d'entiers tels que} @@M5@@, \uline{et} @@M6@@ \uline{le schéma des drapeaux de type} @@M7@@ \uline{de} @@M8@@. \uline{Alors l'homomorphisme canonique} :
@@M9@@
\uline{est un isomorphisme}.

On reprend les notations de VI 4.5. On note donc @@M10@@ les quotients canoniques de l'image inverse de @@M11@@ sur @@M12@@, @@M13@@ le noyau de @@M14@@ (@@M15@@), et @@M16@@ la classe dans @@M17@@ de @@M18@@. Alors les @@M19@@ engendrent @@M20@@ comme @@M21@@-algèbre (VI 4.6). De même, les images des @@M22@@ dans @@M23@@ engendrent @@M24@@ comme algèbre sur @@M25@@ pour tout point @@M26@@ de @@M27@@. Comme @@M28@@ est un schéma de drapeaux sur un corps, c'est un schéma régulier, donc l'homomorphisme @@M29@@ est un isomorphisme, si bien que l'homomorphisme @@M30@@ est surjectif. D'après 1.3, cela entraîne que l'homomorphisme:
@@M31@@
est surjectif.

Pour montrer l'injectivité, on peut se ramener d'abord au cas du fibré en drapeaux complet @@M32@@. En effet, supposons que
@@M33@@
soit injectif. Comme @@M34@@ se factorise par @@M35@@, cet
